% ============================================================
% resume AIAgent Compact.tex — 精简版，源自 src/resume AIAgent.tex
% 改动：各家项目各压成 3 条要点（去掉 STAR 标签、A/R 合并），去掉「个人项目」；
%       新增交个朋友（高级 AI 开发工程师）：电商视频 Agent 工作台、千川为独立项目，AI Studio、片段工场等并入「其他交付」
% 样式与原版共用 templates/style-aiagent.sty
% 编译（仓库根目录）：xelatex -output-directory=output "src/resume AIAgent Compact.tex"

\documentclass[10pt,a4paper]{article}

% --- 加载样式与模板 ---
\input{templates/style-aiagent.sty}
\input{templates/cv.latex}

% --- 局部扩展：四列布局等（调格式直接改数字） ---
\let\cvheadingThree\cvheading% 保存3列版本，项目部分恢复用
% 教育经历专用四列（时间 / 学校 / 学历 / 专业）
\newcommand{\cvheadingEdu}[4]{%
    \noindent
    \makebox[0.32\textwidth][l]{\textbf{#1}}% 时间
    \makebox[0.33\textwidth][l]{\textbf{#2}}% 学校
    \makebox[0.12\textwidth][l]{\textbf{#3}}% 学历
    \makebox[0.23\textwidth][r]{\textbf{#4}}\par% 专业
}
\newcommand{\cveventEdu}[5]{%
    \vspace{0.3em}
    \cvheadingEdu{#1}{#2}{#3}{#4}
    \ifx\relax#5\relax\else
        \vspace{0.06em}
        #5
    \fi
}
% 五列版本：工作经历用（时间 / 公司 / 领域 / 语言 / 岗位）
\newcommand{\cvheadingWork}[5]{%
    \noindent
    \makebox[0.27\textwidth][l]{\textbf{#1}}% 时间
    \makebox[0.17\textwidth][l]{\textbf{#2}}% 公司（图标+名称）
    \makebox[0.08\textwidth][l]{\textbf{#3}}% base地
    \makebox[0.20\textwidth][l]{\textbf{#4}}% 行业领域
    \makebox[0.28\textwidth][r]{\textbf{#5}}\par% 岗位+语言
}
% 工作经历专用 \cvevent（6参：时间/公司/base地/领域/岗位+语言/描述）
\newcommand{\cveventWork}[6]{%
    \vspace{0.2em}
    \cvheadingWork{#1}{#2}{#3}{#4}{#5}
    \ifx\relax#6\relax\else
        \vspace{0.04em}
        #6
    \fi
}
% section 标题旁的技术栈图标（PNG 彩色小图标）
% 技术栈图标：优先 devicon(-original) → simple-icons(-simple) → 手动(-custom)
\newcommand{\ticon}[1]{%
    \IfFileExists{assets/myIcons/#1-original.png}{%
        \raisebox{-0.08cm}{\includegraphics[height=0.42cm]{assets/myIcons/#1-original.png}}%
    }{\IfFileExists{assets/myIcons/#1-simple.png}{%
        \raisebox{-0.08cm}{\includegraphics[height=0.42cm]{assets/myIcons/#1-simple.png}}%
    }{%
        \raisebox{-0.08cm}{\includegraphics[height=0.42cm]{assets/myIcons/#1-custom.png}}%
    }}%
}
\newcommand{\seclogo}[1]{\hspace{0.12em}\ticon{#1}}

% ============================================================
% 个人信息设置
% ============================================================
\name{赵志文}
\lineA{%
    \faPhone\ 185-5613-5826 \enspace|\enspace
    \faEnvelope\ \href{mailto:cnjszzw@qq.com}{cnjszzw@qq.com} \enspace|\enspace
    \faWeixin\ cnjszzw%
}
\lineB{籍贯：江苏扬州 \enspace|\enspace 年龄：27 (1999/10/27) \enspace|\enspace 性别：男}
\lineC{求职意向：软件开发}
\photo{assets/photo.jpg}

% ============================================================
% 文档正文
% ============================================================
\begin{document}

\makecvheader

% ============================================================
% 教育经历
% ============================================================
\vspace{-0.825cm}%  第一个标题向上偏移，利用 header 下方空白
\cvsection{教育经历}

% 切换为4列布局（学校 / 学历 / 专业·学院·地点 / 时间）
\let\cvheading\cvheadingEdu
\let\cvevent\cveventEdu

\cvevent{2022/09 -- 2025/01}{%
    \raisebox{-0.12cm}{\includegraphics[height=0.50cm]{assets/school/Xijing_University_Emblem.png}}%
    \hspace{0.4em}西京学院}{硕士}{计算机技术}{}

\cvevent{2018/09 -- 2022/06}{%
    \raisebox{-0.12cm}{\includegraphics[height=0.50cm]{assets/school/Shandong_Normal_University_seal.png}}%
    \hspace{0.4em}山东师范大学}{本科}{广播电视编导}{}

% ============================================================
% 工作经历
% ============================================================
\cvsection{工作经历}

% 切换为4列布局（时间 / 公司 / 领域 / 岗位）
\let\cvheading\cvheadingWork
\let\cvevent\cveventWork

\cvevent{2026/07 -- 至今 (正职)}{%
    \raisebox{-0.12cm}{\includegraphics[height=0.55cm]{assets/LOGO/befriends.jpeg}}%
    \hspace{0.4em}交个朋友}{杭州}{电商/AIGC}{(Python)高级 AI 开发工程师}{}

\cvevent{2025/04 -- 2026/03 (正职)}{%
    \raisebox{-0.12cm}{\includegraphics[height=0.55cm]{assets/LOGO/上海环创通信科技股份有限公司.png}}%
    \hspace{0.4em}寰创通信}{上海}{物联网(政企B端)}{(Java/Python)全栈开发}{}

\cvevent{2024/02 -- 2024/05 (实习)}{%
    \raisebox{-0.12cm}{\includegraphics[height=0.55cm]{assets/LOGO/美团.png}}%
    \hspace{0.4em}美团}{上海}{电商(B端商家侧)}{(Java)测试开发}{}

\cvevent{2023/09 -- 2024/02 (实习)}{%
    \raisebox{-0.12cm}{\includegraphics[height=0.55cm]{assets/LOGO/哔哩哔哩.png}}%
    \hspace{0.4em}哔哩哔哩}{上海}{内容社区(C端)}{(Python)测试开发}{}

% ============================================================
% 项目经历（按公司分组）
% ============================================================
\let\cvheading\cvheadingThree% 恢复3列版本

\cvsection{交个朋友}[%
    \seclogo{claudecode}\seclogo{cursor}%                                    AI 编程工具
    \seclogo{python}\seclogo{fastapi}\seclogo{pytest}\seclogo{ffmpeg}\seclogo{nodejs}%  Python / 媒体
    \seclogo{vuejs}\seclogo{mysql}\seclogo{redis}\seclogo{docker}\seclogo{git}%         平台与基建
]

\cvproject{项目一：电商视频 Agent 工作台}{}{2026/07 -- 至今}{%
    \cvtech{Python、FastAPI、WebView2、FFmpeg、Django、Vue 3、SQLite}
    \begin{itemize}
        \item \textbf{一键成片（老罗版）}：输入商品链接即可自动采集商品事实、生成卖点口播，套用老罗 IP 的形象与话术，按经典六段结构产出竖屏成片，单批可选商品数与条数；采集、卖点、文案、配音、首帧、合成串成一条可断点续跑的父任务，升级后点「继续原任务」只补缺失片段，已完成的文案、配音与视频全部复用。
        \item \textbf{爆款复刻}：把一条爆款原片拆成「洗片 \thickarrow{} 拆片 \thickarrow{} Hypit 复刻 \thickarrow{} 合成」的主流程 DAG——洗片先按镜头判断有没有人，只对有人镜头换人换声，无人镜头原样保留以省生成费；拆片沿完整镜头切开，配词级 ASR 与动作事件生成结构化片段；复刻环节调用 Hypit（RunningHub H3 / Seedance）按原镜头节奏重拍；成片按原顺序拼回，超长镜头按边界合并提交而不截断。
        \item \textbf{混剪}：以词级 ASR 定位原话与句尾，结合画面动作证据确定可用时段；组装阶段校验商品权益一致性、相邻信息增量与源片段重叠率，保留源区间便于审片追溯；同组素材规则迭代后可用成片由 3 条增至 22 条（单组复测口径）。
    \end{itemize}
}

\cvproject{项目二：巨量千川预警与投放调控平台}{}{2026/07 -- 至今}{%
    \cvtech{Django、Celery、MySQL、Redis、Vue 3、TypeScript、Docker}
    \begin{itemize}
        \item \textbf{多账户预警}：止损、追爆、空耗（高消耗低转化）、退货四类规则扫描投放计划，命中后按接收范围聚合飞书卡片并保存发送进度，已确认发送不重发；单次人工验收覆盖 5 个账户、249 个计划，产生 79 条告警。
        \item \textbf{自动追投}：按素材归因的订单、ROI 与退款核对筛出可追素材，经预算、投放时长、冷却窗口与抖音号黑名单约束后调用官方素材追投接口，在途去重避免重复追投。
        \item \textbf{异步任务与写安全}：任务先持久化再投递，重启从 MySQL 检查点续跑，执行者凭带期限的租约提交；有财务副作用的写入记录写前快照与幂等键、写后强制回读，回包不明先查真实状态，确认失败才停自动写入；账户停用后排队任务重新校验当前状态。
        \item \textbf{批量建计划}：ZIP 素材与 XLSX 参数逐行校验，记录文件摘要与行级结果，成功行跳过、失败行按原输入续跑，已上传视频按记录复用。
    \end{itemize}
}

\cvproject{其他交付}{}{2026/07 -- 至今}{%
    \begin{itemize}
        \item \textbf{AI Studio 工作台与 Local Runner}：可恢复的数字人口播 DAG（幂等键 + SQLite 提交记录 + 服务商任务 ID，状态不明不重投、只按原任务回查），网页到本机的短期凭据交接与受管 Skill 同步（按来源 + 版本校验后原子替换）。
        \item \textbf{片段工场与复刻工作台}：资产审核回执以记录 SHA-256 绑定学习版本，学习内容变更即置 review\_stale；4K 预处理改为共享分析代理，固定 12.13 秒样本 CPU 由 730.43 秒降至 60.38 秒（约 -91.7\%），346 条资产首屏响应由约 6.3 MB 降至 73 KB。
        \item \textbf{视频 Skill 与素材同步}：封装 2 个 Codex Skill（13 套版本化模板、按条预算预占与付费幂等、Seedance 与 MiniMax H3 同输入对照）；开发 Windows 素材同步工具（同卷暂存校验 + 原子落位 + 飞书回写回读 + 一键自检）。
    \end{itemize}
}

\needspace{11\baselineskip}% 上海段整段落到下一页，避免标题卡在页边
\cvsection{上海寰创通信股份有限公司}[%
    \seclogo{claudecode}\seclogo{cursor}%                                  AI 工具
    \seclogo{python}\seclogo{langchain}\seclogo{langgraph}\seclogo{deepseek}\seclogo{mcp}\seclogo{fastapi}% Python/AI 栈
    \seclogo{java}\seclogo{spring}\seclogo{springboot}\seclogo{mybatis}\seclogo{mybatis-plus}\seclogo{mysql}\seclogo{redis}\seclogo{mqtt}% Java 栈
    \seclogo{vuejs}\seclogo{javascript}\seclogo{typescript}\seclogo{cesium}% 前端
    \seclogo{git}\seclogo{docker}\seclogo{jenkins}\seclogo{subversion}%    基础设施
]

\cvproject{项目一：无人机管理云平台}{}{2025/04 -- 2026/03}{%
    \cvtech{Java、SpringBoot、Redis、MySQL、MyBatis(Plus)、WebSocket、MQTT、Vue、JavaScript、TypeScript、Cesium}
    \cvlabeled{项目概述}{在公司物联网平台基础上，利用自组网超远距离（>20km）优势新增无人机接入模块，实现空地一体化指挥调度。}
    \cvlabeled{项目实现}{%
        \begin{itemize}
            \item \textbf{无人机接入与设备管理}：将无人机纳入现有设备管理体系并复用区域权限；Redis 维护在线状态（60s TTL 心跳续期）与离线原因，区分主动下线与异常断连；监听 Spring 启动就绪事件，服务启动后经 MQTT 向硬件逐台下发重启指令，重启后自动恢复在线，无需人工值守。
            \item \textbf{飞行控制与相机系统}：实现一键起飞、指点飞行、环绕飞行、返航、急停及变焦、拍照、录像、云台重置等飞行与相机指令；前期前端直连 MQTT 快速验证，后期将关键相机指令与航线任务迁至后端，并补充限高校验与降落确认，降低超限与误操作风险。
            \item \textbf{无人机直播观看}：基于 Spring DeferredResult 异步确认——发起观看即释放线程，由 ZLM 流注册回调（12s 超时兜底）确认推流就绪后返回播放信息；Redis 按路维护观看用户引用计数，人数归零才停止推流，缺流时重新推流，实现多人共享同一路流。
        \end{itemize}
    }
}

\cvproject{项目二：WVP Copilot --- 基于 LLM 的无人机 AI Agent 智能体}{}{2025/11 -- 2026/03}{%
    \cvtech{Python、LangChain、LangGraph、SSE、MCP、FastAPI、DeepSeek API、Redis}
    \cvlabeled{项目概述}{面向消防应急指挥的无人机 Copilot：把自然语言需求（如“飞到陆家嘴，变焦 7x 拍照，返航”）拆解为指令序列并编排执行；采用 Python（Agent 编排层）+ Java（设备控制层）两层架构，通过 REST 解耦。}
    \cvlabeled{项目实现}{%
        \begin{itemize}
            \item \textbf{异步状态同步与故障兜底}：工具函数内每秒轮询 Redis 任务状态，超时阈值按预估飞行时间动态设定；超时不直接判失败，先查 OSD 位置兜底（< 5m 判到达、无位移判硬件故障、仍在移动则延长等待），中间进度只推前端，LLM 只取最终状态。
            \item \textbf{任务链级联阻断与安全隔离}：Plan-Confirm-Execute 三段式，计划经用户确认卡片审核后逐步执行，任一步返回失败或急停即中断后续；飞行、返航在工具层内置确认钩子；急停走独立 MQTT 通道并同步中断 SSE 连接，不经过 Agent 与 LLM。
            \item \textbf{状态机替代 LLM 时序决策}：用 LangGraph 条件子图建模环绕飞行——enter\_orbit \thickarrow{} orbit\_step（按角度步进拍照，$\ge 360^\circ$ 条件边退出）\thickarrow{} exit\_orbit，拍照时机由图状态转移决定，规避 LLM 推理延迟在 3m/s 下造成的 3-9 米偏移；线性主链路仍用 AgentExecutor。
        \end{itemize}
    }
}

% ============================================================
% 专业技能
% ============================================================
\cvsection{专业技能}

\begin{multicols}{2}
\cvskill{Java 全家桶}{熟练掌握 Java，熟悉 JVM 内存模型与并发编程；熟练使用 SpringBoot、MyBatis（Plus）；了解 SpringCloud 微服务体系（Nacos、OpenFeign、Gateway）}

\cvskill{数据存储与检索}{熟悉 MySQL，有 N+1 查询优化（预加载 + HashMap 索引）、延迟关联、覆盖索引等 SQL 调优实战经验；熟悉 Redis 缓存策略与数据一致性保障；熟悉 Elasticsearch 全文检索（Multi-Match、高亮、多维度排序）}

\cvskill{AI Agent 工程化}{有 LLM Agent 从 0 到 1 的实践经验：熟练使用 LangChain、LangGraph、Function Calling、MCP 协议；掌握 Agent 可靠性设计；熟练使用 Cursor/ClaudeCode}

\cvskill{分布式系统}{了解分布式锁（Redisson）、分布式事务（RocketMQ 事务消息 / TCC / 可靠消息最终一致）、配置中心（Nacos）等常见组件的设计思路与应用场景}

\cvskill{性能优化}{有核心接口性能优化实战经验：千台 Mesh 设备查询从 7-8s 优化至 1s 以内，30w 条日志深分页从 5s 降至 1s 以内（性能提升 5 倍）；善于定位并解决大数据量场景下的查询与存储瓶颈}

\cvskill{工程实践}{熟悉 Logback 日志框架与 Kibana 排障，熟悉 Linux + Docker 基础运维；具备从需求分析到联调交付的全流程能力，有前后端 + 软硬件跨团队协作经验}
\end{multicols}

% ============================================================
% 竞赛/获奖/证书
% ============================================================
\cvsection{竞赛/获奖/证书}

\begin{itemize}
    \item 考研初试368分（专业第一）\enspace|\enspace 2022西京学院一等奖学金（前20\%）\enspace|\enspace 2022传智杯程序设计赛区域赛二等奖（前20\%）
    \item 雅思IELTS（6.5分）\enspace|\enspace 英语六级CET-6（593分）\enspace|\enspace 国际人才英语考试ETIC（优秀等级）\enspace|\enspace 计算机二级证书（MS OFFICE）
\end{itemize}

\end{document}
